\section{Comparisons and rational Lefschetz classes}
\label{sec:preliminaries}

We first put the realizations in Theorem~\ref{thm:main} in a common
notation and formulate its conclusion as a property of cohomology classes.
The essential functorial fact is that this property is preserved by rational
pullback. We also record the criterion that will lift an additional divisor
once a suitable Hodge filtration has been realized.

We work in the category of abelian varieties up to isogeny. Thus
\(\Hom^0(X,Y)=\Hom(X,Y)\otimes_\Z\Q\), and divisors may have rational
coefficients. Pullback by a rational homomorphism is defined on \(H^1\)
and extended to its exterior algebra. If \(Nf\) is an actual homomorphism,
then on \(H^{2r}\)
\begin{equation}\label{eq:rational-pullback}
 f^*=N^{-2r}(Nf)^*.
\end{equation}
This convention is the same in all realizations.

For an abelian variety \(X/\F\), set
\[
 k_\ell=\Q_\ell\quad(\ell\ne p),\qquad k_p=\Cw,
\]
and write \(H_v^{2r}(X)(r)\) for its \(\ell\)-adic cohomology if
\(v=\ell\ne p\), and for crystalline cohomology extended to \(\Cw\)
if \(v=p\). Cup products will usually be written without a cup symbol.
All trace maps and divisor classes have their cycle-class normalizations.
The Betti Tate line \(\Q(r)\) supplies the twists in the other
realizations through comparison. Thus \(\Tr_v\) is the trace on top
cohomology with its top Tate twist and sends the class of a zero-cycle to
its degree. For a good-reduction lift, \(\gamma_{0,\ell}\) denotes
smooth proper specialization of its absolute-Hodge realization, and
\(\gamma_{0,p}\) denotes de Rham--crystalline specialization extended
from \(\Kzero\) to \(\Cw\). We keep first crystalline cohomology
over \(\Kzero\) when constructing filtrations, and specify each
extension of the filtration field separately.

We use the exterior-algebra description of abelian cohomology and
faithfulness of rational homomorphisms on \(H^1\); see
\cite[Theorems~12.5 and~15.1]{MilneAV} for the prime-to-characteristic
realizations. The uniform Weil-cohomology formalism, including
crystalline cohomology, exterior powers, cycle-normalized traces and
the projection formula, is set out in
\cite[Section~1 and Appendix, Proposition~(A.2)]{Milne1999}. Smooth
proper specialization away from \(p\) is
\cite[Theorem~20.4]{MilneLEC}.

\subsection{Comparison of Hodge tensors}

\begin{lemma}\label{lem:comparison}
For an abelian variety with good reduction, the realizations of an
absolute Hodge class are compatible with the de Rham and crystalline
comparison isomorphisms. In particular, the \(p\)-adic crystalline tensor
obtained from its \(p\)-adic \'etale realization agrees, after extension of
scalars, with the de Rham--crystalline realization in
Theorem~\ref{thm:main}. These compatibilities respect tensor operations,
duals, Tate twists, divisor classes, and rational homomorphisms.

The same assertion applies to auxiliary abelian varieties with good
reduction over a finite extension of \(\Kzero\), using a complex
embedding of an algebraic closure of their field of definition.
\end{lemma}

\begin{proof}
For abelian varieties over \(\Qbar\), Deligne's absolute-Hodge theorem
\cite[Main Theorem~2.11]{Deligne1982} and Blasius's comparison theorem
\cite[Theorems~0.3 and~5.3]{Blasius1994} give the assertion. The
compatibility of the crystalline comparison with the de Rham comparison
is part of \cite[\S5.1]{Blasius1994}. The formulation for a complete
mixed-characteristic discrete valuation ring with perfect residue field
is given in \cite[Theorem~5.6.3]{Moonen1998}; the paragraph following
that theorem explains removal of the number-field hypothesis. A finite
extension of \(\Kzero\) is within that scope. The stated functorialities
are those of the comparison isomorphisms.
\end{proof}

An abstract complex embedding used here need not be continuous. To see
that it can be chosen compatibly with \(\Qbar\subset\C\), take a
transcendence basis of \(\overline{\Kzero}/\Qbar\), embed it in a
transcendence basis of \(\C/\Qbar\), and extend algebraically. The
cardinalities permit this construction. No descent of a general
auxiliary lift to \(\Qbar\) is inferred.

\subsection{The common rational Lefschetz space}

Let \(\Lef^r(X)\) be the space of rational polynomial expressions in
divisors of codimension \(r\), modulo numerical equivalence. We use the
following established properties:
\begin{enumerate}[label=(\roman*)]
 \item its cycle map is injective, and
 \(\Lef^r(X)\otimes_\Q k_v\) is the divisor-generated subspace of
 \(H_v^{2r}(X)(r)\);
 \item intersection and degree give a perfect pairing
 \[
   \Lef^r(X)\times\Lef^{\dim X-r}(X)\longrightarrow\Q;
 \]
 \item pushforward by a homomorphism of abelian varieties preserves the
 divisor-generated subspaces in each realization;
 \item hard Lefschetz and its inverse restrict to these subspaces.
\end{enumerate}
These are \cite[Proposition~5.2, Corollaries~5.3 and~5.5,
Theorem~5.9]{Milne1999}, with the author's correction to the proof of
Theorem~5.9. Crystalline cohomology is included. Extension of its
coefficient field to \(\Cw\) preserves these statements. We call the
tuple of realizations of a fixed element of \(\Lef^r(X)\) a
\emph{simultaneous rational Lefschetz class}.

\begin{definition}
A tuple \(\alpha=(\alpha_v)_v\), with
\(\alpha_v\in H_v^{2r}(X)(r)\), has the \emph{Lefschetz pairing
property} if, for every \(\delta\in\Lef^{\dim X-r}(X)\), the numbers
\(\Tr_v(\alpha_v\delta_v)\) are the images of one number in \(\Q\).
\end{definition}

Simultaneous rational algebraic classes have this property, by the
intersection formula. The property is closed under rational linear
combinations. We shall need the following additional closure property.

\begin{lemma}\label{lem:pairing-functorial}
If \(\alpha\) has the Lefschetz pairing property on \(Y\) and
\(f\in\Hom^0(X,Y)\), then \(f^*\alpha\) has that property on \(X\).
\end{lemma}

\begin{proof}
If \(r>\dim Y\), the input class is zero; if \(r>\dim X\),
its pullback is zero. We may therefore suppose that \(f\) is an actual
homomorphism and \(r\leq\min(\dim X,\dim Y)\).
For \(\delta\in\Lef^{\dim X-r}(X)\), the class \(f_*\delta_v\)
lies in the Lefschetz subspace of codimension \(\dim Y-r\) in every
realization. Choose rational bases \(a_i\) of
\(\Lef^{\dim Y-r}(Y)\) and \(b_j\) of \(\Lef^r(Y)\). The matrix
\(M_{ij}=\deg(a_i b_j)\) is an invertible rational matrix. If
\(f_*\delta_v=\sum_i c_{i,v}(a_i)_v\), the projection formula gives
\[
 \sum_i c_{i,v}M_{ij}=\deg(\delta f^*b_j).
\]
The right side is a rational number independent of \(v\). Solving this
one rational system shows that \(c_{i,v}=c_i\in\Q\) for every \(v\).
Thus \(f_*\delta\) is simultaneous rational Lefschetz. Applying the
projection formula once more proves the assertion. For rational \(f\),
use \eqref{eq:rational-pullback}.
\end{proof}

\subsection{Lifting a prescribed rational homomorphism}

The second transport construction will produce an abelian lift with a
chosen Hodge filtration. The following criterion then lifts the algebraic
projector, field action, and divisor already specified on its special fiber.
Its input is a rational algebraic homomorphism, rather than an arbitrary
map between cohomology spaces.

In the next lemma \(R\) is a complete mixed-characteristic discrete
valuation ring finite over \(W(\F)\), with fraction field \(K\).
For an abelian scheme \(\mathcal X/R\), the Hodge filtration is regarded
as a filtration of
\(H^1_{\mathrm{cris}}(X_0/W(\F))[1/p]\otimes_{\Kzero}K\).

\begin{lemma}\label{lem:lift-hom}
Let \(\mathcal X,\mathcal Y\) be abelian schemes over \(R\), with
special fibers \(X_0,Y_0\), and let
\(f_0\in\Hom^0(X_0,Y_0)\). If its contravariant crystalline map
preserves the Hodge filtrations, then \(f_0\) lifts uniquely to a
rational homomorphism \(\mathcal X\to\mathcal Y\).

In particular, a symmetric rational homomorphism \(X_0\to X_0^\vee\)
with this property lifts to a rational divisor class on the generic
fiber. Algebraic identities among finitely many such lifted maps remain
valid.
\end{lemma}

\begin{proof}
The crystalline map already commutes with Frobenius. Full faithfulness
for rational filtered Dieudonn\'e modules gives a rational morphism of
the associated \(p\)-divisible groups; see
\cite{Grothendieck1970} and the perfect-residue-field formulation in
\cite[\S4.2(v)]{IllusiePisa}. Equivalently, crystalline comparison first
gives the rational map on Tate modules. Multiply by an integer \(N\)
that clears both its Tate-lattice denominator and the denominator of
\(f_0\). Tate's extension theorem
\cite[Theorem~4]{Tate1967} extends the resulting integral generic-fiber
map of \(p\)-divisible groups over \(R\). Faithfulness of the
special-fiber Dieudonn\'e functor identifies its reduction with the
map induced by \(Nf_0\). Here equality after inverting \(p\) is
already equality of integral maps, because the integral Dieudonn\'e
modules are torsion free. The integer \(N\) is fixed for the remainder
of the construction.

The specified-map form of Serre--Tate deformation theory now gives a
lift over each \(R/\mathfrak m_R^n\)
\cite[p.~432]{Grothendieck1970}. One can also deduce it from the
fixed-object formulation \cite[Theorem~2.3.3]{IllusiePisa} as follows.
On \(X_0\times Y_0\) form the shear
\(s_0(x,y)=(x,y+Nf_0(x))\). The lifted \(p\)-divisible-group map gives
the same shear on the product \(p\)-divisible group. View the existing
product deformation with two special-fiber markings differing by
\(s_0\). Full faithfulness lifts the shear to an algebraic isomorphism
between these marked deformations; its second component restricted to
\(X\) lifts \(Nf_0\). Uniqueness makes these maps compatible at every
finite order.

Formal existence algebraizes the resulting formal map
\cite[Tag~0A42]{StacksProject}: its source is proper over the complete
noetherian ring \(R\), and its target is separated and of finite type.
Dividing by the fixed \(N\) gives the desired rational lift. Rigidity
gives uniqueness and injectivity of specialization on rational
homomorphisms \cite[Proposition~20.1 and Corollary~20.2]{MilneAV}.

If \(f_0\) is symmetric, injectivity shows that its lift is symmetric.
The usual identification
\[
 \NS(X_K)\otimes\Q
   \simeq\Hom(X_K,X_K^\vee)^{\mathrm{sym}}\otimes\Q
\]
gives the rational divisor class, after a finite extension if
necessary \cite[Proposition~17.2]{MilneAV}. Finally an identity between the maps specializes to the
given identity, so injectivity proves it upstairs as well.
\end{proof}

Only prescribed rational algebraic maps on the special fiber are being
lifted in Lemma~\ref{lem:lift-hom}. An arbitrary \(\Q_p\)-linear map of
\(p\)-divisible groups is not asserted to be a rational algebraic map.
We also use the extension of generic-fiber homomorphisms between abelian
schemes over a discrete valuation ring: their N\'eron mapping property
extends such a homomorphism uniquely over the base
\cite[Definition~5.1 and Example~5.3]{Conrad2003}. Clearing one
rational denominator therefore makes specialization of these maps
well defined.
